\section*{Phụ lục}
\addcontentsline{toc}{section}{Phụ lục}
\renewcommand{\thesubsection}{\Alph{subsection}.}
\setcounter{subsection}{0}

\subsection{Mã Power Query}
\choviet{report-draft}

\subsection{Nhật ký ngoại lệ đầy đủ}
Sinh từ truy vấn \texttt{q\_ExceptionLog} (sheet \texttt{Nhat\_ky\_du\_lieu}); sheet \texttt{Nhat\_ky\_ngoai\_le} của file mẫu đọc lại đúng bảng này. Cột ``Phát hiện ở bước'' và ``Nguồn'' xem trong file Excel.

{\scriptsize
\begin{longtable}{@{}L{0.85cm} L{4.05cm} r L{1.6cm} L{1.6cm} L{2.85cm} L{1.7cm}@{}}
\toprule
Mã & Mô tả & USD & Kỳ & Phân loại & Hướng xử lý & Phụ trách; trạng thái \\
\midrule
\endhead
EX-01 & Flash ghi giá niêm yết, sổ cái ghi giá sau giảm giá: tổng chiết khấu của mọi dòng POS và đơn online & $-$1.834.297,47 & 2025-01 đến 2025-12 & Quy tắc ghi nhận & Trình bày trong cầu nối; đề nghị Flash báo thêm doanh số sau giảm giá & Kinh doanh; Đã giải thích \\
EX-02 & 12 dòng POS trùng khóa ở cửa hàng NX007 (mã POS 74) ngày 29/03/2025: file POS xuất 2 lần, Flash đếm cả bản trùng & $-$12.622,20 & 2025-03 & Lỗi dữ liệu & Bỏ bản xuất sau; chặn xuất lặp ở hệ thống POS & IT; Cần sửa nguồn \\
EX-03 & 124 đơn online đặt trong kỳ rồi bị hủy: Flash có tính, sổ cái không ghi & $-$215.658,92 & 2025-01 đến 2025-12 & Quy tắc ghi nhận & Đề nghị Flash loại đơn hủy hoặc báo riêng & Kinh doanh; Đã giải thích \\
EX-04 & Đơn đặt và giao khác tháng trong năm: tháng này chuyển đi, tháng sau nhận lại; lớn nhất 2025-03 (281.787,98) & 0,00 & 2025-01 đến 2025-12 & Quy tắc ghi nhận & Không sửa; ghi chú trên Flash rằng doanh thu tính theo ngày đặt & Kinh doanh; Đã giải thích \\
EX-05 & 77 đơn đặt trước kỳ, giao trong kỳ: sổ cái có, Flash của kỳ không có & 122.107,89 & 2025-01 & Quy tắc ghi nhận & Không sửa; trình bày riêng khi so cả năm & Kinh doanh; Đã giải thích \\
EX-06 & 105 đơn đặt trong kỳ, giao sau kỳ: Flash có, sổ cái kỳ sau mới ghi & $-$174.705,79 & 2025-12 & Quy tắc ghi nhận & Theo dõi: sổ cái kỳ sau phải ghi nhận đủ số này & Tài chính; Đã giải thích \\
EX-07 & 65 dòng POS bán cuối tháng 2025-05 vào ngày nghỉ, ghi sổ tháng 2025-06: 41.464,78 USD chuyển tháng, tự bù trừ & 0,00 & 2025-05, 2025-06 & Quy tắc ghi nhận & Không sửa; tự bù trừ ở tháng sau & Tài chính; Đã giải thích \\
EX-08 & 101 dòng POS bán cuối tháng 2025-08 vào ngày nghỉ, ghi sổ tháng 2025-09: 75.240,30 USD chuyển tháng, tự bù trừ & 0,00 & 2025-08, 2025-09 & Quy tắc ghi nhận & Không sửa; tự bù trừ ở tháng sau & Tài chính; Đã giải thích \\
EX-09 & 101 dòng POS bán cuối tháng 2025-11 vào ngày nghỉ, ghi sổ tháng 2025-12: 80.588,85 USD chuyển tháng, tự bù trừ & 0,00 & 2025-11, 2025-12 & Quy tắc ghi nhận & Không sửa; tự bù trừ ở tháng sau & Tài chính; Đã giải thích \\
EX-10 & NX057 ngày 21/11/2025: POS có 7.227,56 USD nhưng sổ cái không có bút toán & $-$7.227,56 & 2025-11 & Lỗi dữ liệu & So ô cửa hàng × ngày trước khi khóa sổ; IT kiểm tra giao diện chuyển dữ liệu POS sang sổ cái & Tài chính, IT; Đã sửa kỳ sau \\
EX-11 & NX057 ngày 02/12/2025: sổ cái có bút toán JE20251206231 (POS\_CATCHUP) nhưng POS không có dòng nào & 7.227,56 & 2025-12 & Lỗi dữ liệu & So ô cửa hàng × ngày trước khi khóa sổ; IT kiểm tra giao diện chuyển dữ liệu POS sang sổ cái & Tài chính, IT; Đã sửa kỳ sau \\
EX-12 & 169 đơn bị trả, sổ cái ghi giảm vào ngày nhận hàng trả; Flash không trừ (tỷ giá ngày đặt) & $-$349.148,23 & 2025-01 đến 2025-12 & Quy tắc ghi nhận & Đề nghị tách báo cáo hàng trả bên cạnh Flash & Kinh doanh; Đã giải thích \\
EX-13 & 12 bút toán tay MJ-SHIP ghi phí vận chuyển vào doanh thu online; Flash không có khoản này & 94.766,95 & 2025-01 đến 2025-12 & Quy tắc ghi nhận & Trình bày riêng trong cầu nối & Tài chính; Đã giải thích \\
EX-14 & Bút toán tay MJ-RCL-2506 chuyển 3.480,00 USD doanh thu của NX019 sang kênh online; toàn công ty không đổi & 0,00 & 2025-06 & Quy tắc ghi nhận & Trình bày ở bảng cầu nối theo kênh & Tài chính; Đã giải thích \\
EX-15 & Flash đổi theo tỷ giá ngày đặt / ngày bán, sổ cái theo tỷ giá cuối tháng ghi sổ & $-$57.598,01 & 2025-01 đến 2025-12 & Quy tắc ghi nhận & Trình bày riêng một dòng; không sửa & Tài chính; Đã giải thích \\
EX-16 & Flash làm tròn từng dòng đến USD, sổ cái làm tròn từng bút toán đến 2 chữ số & 1,01 & 2025-01 đến 2025-12 & Sai số chấp nhận & Ghi một dòng, không đào thêm & Tài chính; Đã giải thích \\
EX-17 & Bút toán JE20250904892 (MJ-ADJ-2509) ghi 1.250,00 USD doanh thu, không gắn kênh, không có giấy tờ làm căn cứ & 1.250,00 & 2025-09 & Chưa giải thích được & Yêu cầu phòng Tài chính gửi chứng từ; không gán vào khoản khác & Tài chính; Chờ phản hồi \\
\midrule
& \textbf{Tổng (= chênh lệch sổ cái $-$ Flash cả năm)} & \textbf{$-$2.425.904,77} & & & & \\
\bottomrule
\end{longtable}}


\subsection{Bảng tra nhãn nguồn số liệu}

Mỗi kết luận trong báo cáo ghi kèm nhãn cho biết con số lấy từ bảng nào. Chữ viết tắt trong nhãn: QC (Quality Control, kiểm soát chất lượng), RC (Reconciliation, đối soát), BR (Bridge, cầu nối), EX (Exception, ngoại lệ), CH (cửa hàng), RET (Return, hàng trả).

\captab{Nhãn nguồn số liệu}\nopagebreak[4]
\begin{longtable}{@{}L{3.0cm} L{6.0cm} L{6.0cm}@{}}
\toprule
Nhãn & Bảng & Nơi xem trong file Excel Power Query \\
\midrule
\endhead
\nguon{QC} & Kiểm tra tổng, 5 bẫy, chất lượng dữ liệu & Sheet \texttt{Kiem\_tra}, \texttt{Bay\_du\_lieu}, \texttt{Chat\_luong\_du\_lieu}; truy vấn \texttt{q\_Check}, \texttt{q\_Traps}, \texttt{q\_DataQuality} \\
\nguon{RC-CH} & Đối soát cửa hàng: theo ngày ghi sổ (sheet \texttt{Doi\_soat\_CH\_ngay}), so với ngày bán (\texttt{So\_ngay\_ban\_ngay\_ghi\_so}), theo tháng có cờ vượt ngưỡng (\texttt{Doi\_soat\_CH\_thang}), lệch ngày ghi sổ (\texttt{Lech\_ngay\_ghi\_so}) & Truy vấn \texttt{q\_RC\_Store}, \texttt{q\_RC\_Store\_DayCompare}, \texttt{q\_RC\_Store\_Month}, \texttt{q\_RC\_Store\_Timing} \\
\nguon{RC-WEB} & Đối soát online: theo ngày giao $\times$ loại tiền (sheet \texttt{Doi\_soat\_online}), theo tháng có cờ vượt ngưỡng (\texttt{Doi\_soat\_online\_thang}), lệch thời điểm online (\texttt{Lech\_thoi\_diem\_online}); bảng 5 lý do (\texttt{Nam\_ly\_do}) và danh mục khoản, nhóm (\texttt{Danh\_muc\_khoan}) & Truy vấn \texttt{q\_RC\_Web}, \texttt{q\_RC\_Web\_Month}, \texttt{q\_Web\_Timing}, \texttt{q\_FiveReasons}, \texttt{q\_BridgeItems} \\
\nguon{RC-RET} & Đối soát hàng trả $\times$ ngày trả $\times$ loại tiền (sheet \texttt{Doi\_soat\_hang\_tra}) & Truy vấn \texttt{q\_RC\_Return} \\
\nguon{BR-Tổng} & Bảng cầu nối toàn công ty (dòng Flash, sổ cái: sheet \texttt{Chenh\_lech\_thang}) & Truy vấn \texttt{q\_BR\_Long}, cột \texttt{bridge} = COMPANY \\
\nguon{BR-CH},\newline \nguon{BR-WEB} & Bảng cầu nối cửa hàng, online (dòng Flash, sổ cái: sheet \texttt{Chenh\_lech\_kenh}) & Truy vấn \texttt{q\_BR\_Long}, cột \texttt{bridge} = STORE, WEB \\
\nguon{EX-xx} & Dòng xx của nhật ký ngoại lệ & Truy vấn \texttt{q\_ExceptionLog} \\
\bottomrule
\end{longtable}

\subsection{Ba phiên bản số làm tròn}

Đề có ba nguồn cho cùng một bảng cầu nối cả năm, khác nhau vài USD ở ba khoản (Bảng~\ref{tab:lamtron}). Tổng các khoản luôn đúng bằng chênh lệch; khác nhau chỉ ở chỗ phần làm tròn nằm ở dòng nào.

\begin{table}[H]
\captab{Ba phiên bản của ba khoản cầu nối cả năm, USD}\label{tab:lamtron}
\centering
\begin{tabular}{@{}L{7.2cm} r r r@{}}
\toprule
Nguồn & Chiết khấu & Tỷ giá & Làm tròn \\
\midrule
Nhóm: đổi tiền từ nội tệ $\div$ tỷ giá, không làm tròn từng dòng (định nghĩa slide 34) & \SoBRChietKhau{} & \SoBRTyGia{} & +\SoBRLamTron{} \\
Slide 32 & $-$1.834.297 & $-$57.598 & +1,01 \\
Dashboard mẫu của mentor: cộng các cột USD đã làm tròn sẵn từng dòng & $-$1.834.298,38 & $-$57.598,50 & +2,33 \\
\bottomrule
\end{tabular}
\end{table}

Lý do nhóm chọn cách đầu: dòng làm tròn khi đó chỉ gồm phần Flash làm tròn từng dòng đến USD (khoảng +0,95) và phần sổ cái làm tròn từng bút toán đến 2 chữ số (khoảng +0,05), đúng nghĩa ``làm tròn'', và mỗi tháng không quá \SoBRLamTronMax{} USD (ngưỡng 5 USD). Số của nhóm ở bảng này lấy từ bảng cầu nối Power Query \nguon{BR-Tổng}; bảng tính kiểm tra viết bằng Python ra cùng số.

\subsection{So sánh khi đổi thứ tự các bước cầu nối}
Cùng một khoản lệch có thể ra số khác nhau tùy bước nào tính trước, dù tổng các khoản không đổi (Bảng~\ref{tab:doithutu}). Ví dụ dòng POS trùng tính trước chiết khấu thì thành \SoDoiTrungMoi{} USD (theo giá trước giảm giá); hàng trả tính sau tỷ giá thì thành \SoDoiTraMoi{} USD (theo tỷ giá ngày trả), còn dòng tỷ giá thành \SoDoiTyGiaMoi{}. Vì vậy thứ tự các bước được chốt trước khi tính và ghi rõ trong quy ước. Nhóm chọn tính chiết khấu trước mọi khoản khác để các khoản sau đều đo theo giá thực bán, và đổi mọi khoản theo tỷ giá ngày đặt để toàn bộ phần do tỷ giá nằm ở một dòng.

\begin{table}[H]
\captab{Số của hai khoản khi đổi thứ tự các bước, USD cả năm toàn công ty \nguon{BR-Tổng}}\label{tab:doithutu}
\centering
{\small
\begin{tabular}{@{}l r r r@{}}
\toprule
Cách đổi thứ tự và khoản bị ảnh hưởng & Thứ tự đã chốt & Khi đổi & Thay đổi \\
\midrule
\multicolumn{4}{@{}l}{\textit{1. Tính dòng POS trùng trước chiết khấu (theo giá trước giảm giá)}} \\
\quad Dòng POS bị trùng & $-$12.622,20 & $-$13.695,13 & $-$1.072,93 \\
\quad Chiết khấu & $-$1.834.297,47 & $-$1.833.224,54 & 1.072,93 \\
\midrule
\multicolumn{4}{@{}l}{\textit{2. Tính đơn hủy trước chiết khấu (theo giá trước giảm giá)}} \\
\quad Đơn online bị hủy & $-$215.658,92 & $-$230.286,39 & $-$14.627,47 \\
\quad Chiết khấu & $-$1.834.297,47 & $-$1.819.670,00 & 14.627,47 \\
\midrule
\multicolumn{4}{@{}l}{\textit{3. Tính tỷ giá trước hàng trả (hàng trả theo tỷ giá ngày trả)}} \\
\quad Hàng trả lại & $-$349.148,23 & $-$349.652,76 & $-$504,53 \\
\quad Tỷ giá & $-$57.598,01 & $-$57.093,48 & 504,53 \\
\bottomrule
\end{tabular}}

\end{table}

\subsection{Khai báo dùng AI}

Nhóm dùng trợ lý AI (trí tuệ nhân tạo) Claude để viết nháp code Power Query, DAX, script kiểm tra và một phần lời văn. Mỗi phần việc có một file plan trong thư mục \texttt{plans/} ghi rõ: nhờ AI làm gì, phần nào do AI viết ra, và nhóm kiểm tra lại bằng cách nào (so với bảng tổng kiểm soát của đề, với slide, với một cách tính độc lập). Bảng tổng hợp đầy đủ:
\choviet{report-draft}
